% Complete Section 3.1, original PDF 31-32 / printed 21-22.
% block: unbounded-definition-heading page:31 type:section
\subsection{{Unbounded operators}}\label{sec:unbounded-operators}

% block: unbounded-domain-definition pages:31,32 type:ordinary-paragraph
{When considering unbounded operators, one must allow for operators whose domain may not be the entire Hilbert space but only a subspace. Thus, by an \emph{unbounded operator} $T$ from a Hilbert space $X$ to a Hilbert space $Y$, we mean a linear operator mapping a subspace $D(T)$ of $X$, called the domain of $T$, into $Y$. The operator $T$ is not necessarily bounded and the domain $D(T)$ is not necessarily a closed subspace. Such operators are called unbounded operators on Hilbert space, even though the case of a bounded linear operator defined on all of $X$ is included. The more accurate term \emph{not-necessarily-everywhere-defined, not-necessarily-bounded linear operators} would be a bit unwieldy! We shall almost always be concerned with \emph{densely defined} operators, meaning that the domain $D(T)$ is a dense subspace of $X$. A typical example to keep in mind is $X=L^2(\Omega)$ for some domain $\Omega\subset\mathbb R^n$, $Y=L^2(\Omega;\mathbb R^n)$, the domain $D(T)$ is the Sobolev space $H^1(\Omega)$, which is a dense subspace of $L^2(\Omega)$, and the operator $T$ is the gradient. That is, $T\phi=\operatorname{grad}\phi\in L^2(\Omega;\mathbb R^n)$ for all $\phi\in H^1(\Omega)$. A variation is the operator $\widetilde T$ with $D(\widetilde T)=\mathring H^1(\Omega)$, which is also dense in $L^2$, again with $\widetilde T\phi=\operatorname{grad}\phi$, but with $\widetilde T\phi$ only defined for $\phi\in\mathring H^1$. Other possibilities would be to define a gradient operator whose domain consists of $C^1$ functions on $\overline\Omega$, or $C^\infty$ functions, or $C^\infty$ functions with compact support inside $\Omega$. Although we might refer to any one of these operators as the gradient, they are different, and it is important to always be clear about the domain. If we say that two unbounded operators $S$ and $T$ are equal, we mean both that $D(S)=D(T)$ and that the actions of the operators coincide.}

% block: unbounded-sum-composition page:32 type:ordinary-paragraph
{Two unbounded operators $S$ and $T$ from $X$ to $Y$ can be added to give an operator $S+T$, with domain $D(S+T)=D(S)\cap D(T)$. If $S$ is an unbounded operator $X\to Y$ and $T$ an unbounded operator $Y\to Z$, then we can similarly define the composition $T\circ S$ with domain $\{\,v\in D(S)\mid Sv\in D(T)\,\}$. One must be careful, since the sum or composition of densely defined operators need not be densely defined.}

% block: null-range-graph-intro page:32 type:ordinary-paragraph
{The \emph{null space} of an unbounded operator $T$, the \emph{range} of $T$, and the \emph{graph} of $T$ are each defined in the obvious way:}

% block: null-range-graph-definitions page:32 type:equation
\[
\begin{aligned}
\mathcal N(T)&=\{\,x\in D(T)\mid Tx=0\,\}\subset X,\\
\mathcal R(T)&=\{\,Tx\mid x\in D(T)\,\}\subset Y,\\
\Gamma(T)&=\{\,(x,Tx)\mid x\in D(T)\,\}\subset X\times Y.
\end{aligned}
\]

% block: graph-norm-intro page:32 type:ordinary-paragraph
{The function $v\mapsto(v,Tv)$ is a one-to-one mapping from the domain of $T$ onto its graph. Applying the norm in $X\times Y$ to $(v,Tv)$ defines a norm on $D(T)$. This norm, called the \emph{graph norm} on $D(T)$, is given by}

% block: graph-norm-equation page:32 type:equation
\[
\|v\|_{D(T)}^2:=\|v\|_X^2+\|Tv\|_Y^2,\qquad v\in D(T),
\]

% block: graph-inner-product-intro page:32 type:ordinary-paragraph
{and is associated to the \emph{graph inner product}}

% block: graph-inner-product-equation page:32 type:equation
\[
\langle v,w\rangle_{D(T)}:=\langle v,w\rangle_X+\langle Tv,Tw\rangle_Y,\qquad v,w\in D(T).
\]

% block: graph-norm-strength page:32 type:ordinary-paragraph
{Obviously the graph norm is stronger than the restriction to $T(D)$ of the $X$ norm, and commonly they are not equivalent. For example, in the example given above of the gradient operator, the graph norm is the usual $H^1$ norm, which is strictly stronger than the $L^2$ norm. We always have}

% block: graph-boundedness-equation page:32 type:equation
\[
\|Tv\|_Y\leq\|v\|_{D(T)},\qquad v\in D(T),
\]

% block: graph-norm-completeness page:32 type:ordinary-paragraph
{so the operator $T$ is a bounded operator from $D(T)$ to $Y$ when the former is endowed with the graph norm. The graph norm makes $D(T)$ an inner product space, or pre-Hilbert space, but it may not be complete.}%
% block: graph-norm-footnote page:32 type:footnote
\footnote{{We could have instead defined the graph norm by $\|v\|_{D(T)}=\|v\|_X+\|Tv\|_Y$ without the squares. This definition gives an equivalent norm, but it would no longer be associated to an inner product.}}
